\chapter{Referencias}


[1] G. Bledt et al., “MIT Cheetah 3: Design and Control of a Robust, Dynamic Quadruped Robot,” in Proc. IEEE/RSJ Int. Conf. Intelligent Robots and Systems (IROS), Madrid, Spain, Oct. 2018, pp. 2245–2252.

[2] C. D. Bellicoso, F. Jenelten, C. Gehring, and M. Hutter, “Dynamic locomotion through online nonlinear motion optimization for quadrupedal robots,” IEEE Robotics and Automation Letters, vol. 3, no. 3, pp. 2261–2268, Jul. 2018, doi: 10.1109/LRA.2018.2794620.

[3] A. Zamani, M. Khorram, and S. A. A. Moosavian, “Dynamics and Stable Gait Planning of a Quadruped Robot,” in Proc. 11th Int. Conf. Control, Automation and Systems (ICCAS), Gyeonggi-do, Korea, Oct. 2011, pp. 146–151.

[4] M. Hutter et al., “ANYmal: A Highly Mobile and Dynamic Quadrupedal Robot,” in Proc. IEEE/RSJ Int. Conf. Intelligent Robots and Systems (IROS), 2016, pp. 38–44, doi: 10.1109/IROS.2016.7758092.

[5] R. Singh and R. Kotecha, “Quadruped Robot Gait Planning for Enhanced Locomotion and Stability,” arXiv preprint, arXiv:2505.02414, May 2025. [Online]. Available: https://arxiv.org/abs/2505.02414

[6] N. Hafner, C. Liu, C. Ishi, and H. Ishiguro, “Quadrupedal Spine Control Strategies: Exploring Correlations Between System Dynamic Responses and Human Perspectives,” arXiv preprint, arXiv:2505.02414, 2025.

[7] D. Kim et al., “Dynamic Locomotion for Passive-Ankle Biped Robots and Humanoids Using Whole-Body Locomotion Control,” Journal Title, vol. XX, no. X, pp. 1–18, 2019, doi: 10.1177/ToBeAssigned.

[8] S. Gilroy et al., “Autonomous Navigation for Quadrupedal Robots with Optimized Jumping through Constrained Obstacles,” arXiv preprint, 2023. [Online]. Available: https://arxiv.org/abs/

[9] J. Di Carlo et al., “Dynamic Locomotion in the MIT Cheetah 3 Through Convex Model-Predictive Control,” in Proc. IEEE Int. Conf. Robotics and Automation (ICRA), 2018, pp. 1–5.

[10] Z. Fu, A. Kumar, J. Malik, and D. Pathak, “Minimizing Energy Consumption Leads to the Emergence of Gaits in Legged Robots,” arXiv preprint, arXiv:2501.xxxx, 2025.

[11] N. Heess et al., “Emergence of Locomotion Behaviours in Rich Environments,” arXiv preprint, arXiv:1707.02286, 2017.

[12] Y. Tassa et al., “DeepMind Control Suite,” arXiv preprint, arXiv:1801.00690, Jan. 2018.

[13] M. Neunert et al., “Whole-Body Nonlinear Model Predictive Control Through Contacts for Quadrupeds,” in Proc. IEEE Int. Conf. Robotics and Automation (ICRA), 2017, pp. 1–9.

[14] “Design, motions, capabilities, and applications of quadruped robots: a comprehensive review,” *Frontiers in Mechanical Engineering*, 2024. :contentReference[oaicite:3]{index=3}

[15] H. Kim and S. Park, “Locomotion stability control of quadruped robot using the combination of ankle, hip, and step strategies and fuzzy logic controller,” 3rd Int. Conf. Robotics, Automation and AI, 2023, pp. 148–155. :contentReference[oaicite:4]{index=4}

[16] “Smart quadruped robotics: a systematic review of design, control, sensing and perception,” *Advanced Robotics*, vol. 39, no. 1, 2024. :contentReference[oaicite:5]{index=5}

[17] “FROBT paper,” Field and Service Robotics, 2007.

[18] “Biomimetic quadruped research,” Biomimetics, 2020.

[19] “Quadruped robot locomotion research,” arXiv, 2018.

[20] “Legged robot control,” arXiv, 2021.

[21] “Quadruped locomotion research,” arXiv, 2025.

[22] “Quadruped robotics research,” arXiv, 2025.

[23] “Quadruped locomotion learning,” arXiv, 2022.

[24] “Sensors article,” Sensors Journal, 2021.

[25] “Applied Sciences article,” Applied Sciences, 2021.

[26] S. Author et al., “Research article,” PeerJ Computer Science, 2021.

[27] “Applied Sciences article,” Applied Sciences, 2023.

[28] “Quadruped robotics paper,” 2020.

[29] “Coupling of Body Mechanics and Control for Robust and Efficient Quadrupedal Gait Generation,” ResearchGate, 2023.

[30] “MASQ: Multi-Agent Reinforcement Learning for Single Quadruped Robot Locomotion,” ResearchGate, 2024.

[31] “Quadruped locomotion study,” IJCCCE, 2020.

[32] “Quadruped locomotion paper,” arXiv, 2024.

[33] “Applied Sciences article,” Applied Sciences, 2021.

[34] “Paper 4926,” Congresso Brasileiro de Automática, 2022.

[35] “Research article,” Wiley Online Library, 2019.

[36] “Quadruped locomotion research,” ScienceDirect, 2020.

[37] G. DeSouza, “ICRA paper,” IEEE International Conference on Robotics and Automation, 2007.

[38] “Comprehensive summary of the Institute for Human and Machine Cognition's experience with LittleDog,” 2007.

[39] “Quadruped robotics chapter,” Springer, 2010.

[40] P. Stone et al., “ICRA paper,” IEEE International Conference on Robotics and Automation, 2004.

[41] “Quadruped locomotion optimization,” arXiv, 2019.

[42] “Multiobjective optimization of a quadruped robot locomotion,” 2021.

[43] “Quadruped locomotion control,” arXiv, 2021.

[44] “Quadruped robot locomotion research,” arXiv, 2021.

[45] “Robotics research article,” Nature Machine Intelligence, 2025.

[46] “ETH Zurich robotics research,” ETH Zurich, 2023.

[47] “Biomimetic robotics article,” Biomimetics, 2021.

























